Always-on personal agents increasingly run locally, keep memory across days and work unattended. We
study two questions in that setting with Jig, an open-source local agent with an isolated action reviewer
(the Sentinel), on a single consumer GPU: how memory should be managed over days, and how safe the agent
is when nobody is watching. We release a living benchmark that drives the real agent, real local models
and real local services, with full provenance for every result (the served model file, its quantisation,
the server build, the hardware, the seeds and the exact commit), a compute policy suited to a shared
personal machine, and a one-command path for others to contribute results. For memory, we compare five
strategies on LongMemEval---bounded full context, a rolling summary, retrieval over stored turns, and
agent-curated memory with periodic consolidation that either keeps or forgets superseded facts---run
through the real agent loop and memory store, reporting accuracy alongside every model token spent. For
safety, we take a deliberately conservative stance: rather than author any attack content, we measure
Jig's existing defence layers (core and custom rules, the Sentinel under several model choices, read-only
research mode and the approval queue) against \emph{published} agent prompt-injection benchmarks run
unchanged through a Jig adapter. This version reports the harness, the verified literature and the
experimental design; the memory pilot and the published-benchmark safety runs are queued in the overnight
compute window and their tables are generated automatically as they finish. All quantitative results will
be labelled pilots with wide intervals until the full runs complete.
